\documentclass[tikz,border=10pt]{standalone}
\usepackage[UTF8]{ctex}
\usetikzlibrary{arrows.meta,backgrounds,fit,positioning,shadows.blur,shapes.geometric}

\definecolor{ink}{HTML}{1F2937}
\definecolor{muted}{HTML}{5B6472}
\definecolor{lineblue}{HTML}{2563EB}
\definecolor{frontendfill}{HTML}{EFF6FF}
\definecolor{backendfill}{HTML}{F0FDF4}
\definecolor{datafill}{HTML}{FFF7ED}
\definecolor{actorfill}{HTML}{F8FAFC}
\definecolor{groupborder}{HTML}{CBD5E1}

\tikzset{
  font=\sffamily,
  >={Latex[length=3mm,width=2mm]},
  actor/.style={
    draw=ink, rounded corners=2pt, fill=actorfill,
    minimum width=28mm, minimum height=13mm,
    align=center, line width=.8pt, font=\bfseries\sffamily
  },
  component/.style={
    draw=ink, rounded corners=2pt, fill=white,
    minimum width=45mm, minimum height=18mm,
    text width=41mm, align=center, inner sep=4pt,
    line width=.8pt, blur shadow={shadow opacity=.10}
  },
  datastore/.style={
    cylinder, shape border rotate=90, aspect=.27,
    draw=ink, fill=white, minimum width=39mm, minimum height=18mm,
    align=center, line width=.8pt
  },
  dependency/.style={-Latex, draw=muted, line width=.8pt},
  request/.style={-Latex, draw=lineblue, line width=1.1pt},
  labeltext/.style={font=\scriptsize\sffamily, fill=white, inner sep=1.5pt, text=muted},
  group/.style={draw=groupborder, rounded corners=5pt, line width=.8pt, inner sep=7mm},
  grouptitle/.style={font=\bfseries\small\sffamily, text=ink, fill=white, inner sep=2pt},
  note/.style={draw=groupborder, rounded corners=2pt, fill=actorfill,
    text width=202mm, align=left, inner sep=5pt, font=\small\sffamily, text=ink}
}

% 在组件框右上角绘制 UML component 标识。
\newcommand{\componentglyph}[1]{%
  \begin{scope}[shift={(#1.north east)},xshift=-2.5mm,yshift=-1.5mm,draw=ink,line width=.55pt]
    \draw (-3mm,-1.7mm) rectangle (0,1.7mm);
    \draw (-4.2mm,0.6mm) rectangle (-2.5mm,1.55mm);
    \draw (-4.2mm,-1.55mm) rectangle (-2.5mm,-0.6mm);
  \end{scope}%
}

\begin{document}
\begin{tikzpicture}[node distance=8mm and 12mm]

  \node[font=\bfseries\LARGE\sffamily, text=ink] (title)
    {UC09 举报与管理员复审——UML 组件图};
  \node[below=1.5mm of title, font=\small\sffamily, text=muted] (subtitle)
    {展示举报提交、复审查询以及恢复公开／转为私密／删除处置所依赖的系统组件};

  % 参与者
  \node[actor, below left=17mm and 92mm of subtitle] (reporter) {注册用户\\（举报人）};
  \node[actor, below=30mm of reporter] (admin) {管理员\\（role=2）};

  % 前端组件
  \node[component, right=16mm of reporter] (videoView)
    {\textbf{VideoDetailView}\\\footnotesize 视频详情与举报入口};
  \node[component, right=16mm of admin, yshift=9mm] (reportView)
    {\textbf{AdminReportView}\\\footnotesize 举报复审列表};
  \node[component, below=6mm of reportView] (previewView)
    {\textbf{\footnotesize AdminVideoPreviewView}\\\footnotesize 视频预览与处置};
  \node[component, right=13mm of reportView, yshift=10mm] (axios)
    {\textbf{request.js / Axios}\\\footnotesize JWT 与统一响应处理};

  % 后端组件
  \node[component, right=15mm of axios, yshift=12mm] (security)
    {\textbf{\footnotesize SecurityConfig + JWT Filter}\\\footnotesize 身份认证与访问控制};
  \node[component, right=12mm of security, yshift=9mm] (videoController)
    {\textbf{VideoController}\\\footnotesize POST /video/report};
  \node[component, below=10mm of videoController] (adminController)
    {\textbf{\footnotesize AdminVideoController}\\\footnotesize /admin/video/*};
  \node[component, right=12mm of videoController] (videoService)
    {\textbf{VideoService}\\\footnotesize 举报校验、计数与状态流转};
  \node[component, below=10mm of videoService] (adminService)
    {\textbf{\footnotesize AdminVideoService}\\\footnotesize 查询、通过、驳回与删除};
  \node[component, below=8mm of adminService] (notificationService)
    {\textbf{\footnotesize NotificationService}\\\footnotesize 生成处置结果通知};

  % 数据访问与基础设施组件
  \node[component, right=12mm of videoService, yshift=-9mm] (mapperLayer)
    {\textbf{MyBatis-Plus Mapper 层}\\\footnotesize VideoReportMapper · VideoMapper};
  \node[datastore, right=13mm of mapperLayer] (mysql)
    {\textbf{MySQL 8}\\\footnotesize video\_reports / videos\\notifications};
  \node[component, below=13mm of mysql] (uploads)
    {\textbf{本地 uploads/}\\\footnotesize 视频与封面文件};

  % 业务请求与组件依赖
  \draw[request] (reporter) -- node[labeltext,above]{提交举报} (videoView);
  \draw[request] (admin) -- node[labeltext,above]{进入复审} (reportView);
  \draw[request] (admin.east) to[out=0,in=180] node[labeltext,below]{预览／处置} (previewView.west);
  \draw[request] (videoView) -- node[labeltext,above]{POST /video/report} (axios);
  \draw[request] (reportView) -- (axios);
  \draw[request] (previewView.east) to[out=0,in=-120] (axios.south);
  \draw[request] (axios) -- node[labeltext,above]{HTTP + JWT} (security);

  \draw[request] (security) -- node[labeltext,above]{举报请求} (videoController);
  \draw[request] (security.south) to[out=-90,in=180]
    node[labeltext,below]{管理员接口} (adminController.west);
  \draw[dependency] (videoController) -- node[labeltext,above]{调用} (videoService);
  \draw[dependency] (adminController) -- node[labeltext,above]{调用} (adminService);

  \draw[dependency] (videoService.east) to[out=0,in=150]
    node[labeltext,above]{新增举报／更新状态} (mapperLayer.north west);
  \draw[dependency] (adminService.east) to[out=0,in=-150]
    node[labeltext,below]{查询明细／执行处置} (mapperLayer.south west);
  \draw[dependency] (adminService.south) -- node[labeltext,right]{处置通知} (notificationService.north);

  \draw[dependency] (mapperLayer) -- node[labeltext,above]{SQL} (mysql);
  \draw[dependency] (notificationService.east) to[out=0,in=-145] (mysql.south west);
  \draw[dependency] (adminService.south east) to[out=-35,in=180]
    node[labeltext,below]{删除媒体文件} (uploads.west);

  % UML component glyphs
  \foreach \n in {videoView,reportView,previewView,axios,security,videoController,
                   adminController,videoService,adminService,notificationService,
                   mapperLayer,uploads}{\componentglyph{\n}}

  % 分层边界放到背景，避免遮盖组件。
  \begin{scope}[on background layer]
    \node[group, fill=frontendfill, fit=(videoView)(reportView)(previewView)(axios)] (frontendGroup) {};
    \node[group, fill=backendfill, fit=(security)(videoController)(adminController)
      (videoService)(adminService)(notificationService)] (backendGroup) {};
    \node[group, fill=datafill, fit=(mapperLayer)(mysql)(uploads)] (dataGroup) {};
  \end{scope}
  \node[grouptitle, anchor=south west] at ([xshift=2mm,yshift=1mm]frontendGroup.north west)
    {Vue 3 前端};
  \node[grouptitle, anchor=south west] at ([xshift=2mm,yshift=1mm]backendGroup.north west)
    {Spring Boot 后端};
  \node[grouptitle, anchor=south west] at ([xshift=2mm,yshift=1mm]dataGroup.north west)
    {数据与文件基础设施};

  \node[note, text width=365mm, below=12mm of previewView, anchor=north west] (stateNote) at (reporter.south west |- previewView.south) {
    \textbf{核心状态规则：}\quad
    已发布（status=1）在举报数达到阈值后进入举报复审（status=2）；管理员可恢复公开（status=1）、
    转为私密（status=3）或删除视频记录及其媒体文件。所有管理接口必须通过 JWT 与 role=2 校验。
  };

\end{tikzpicture}
\end{document}
